%%%%%%%%%%%%%%%%%%%%%%%%%%%%%%%%%%%%%%%%%%%%%%%%%%%%%%%%%%%%%%%%%%%%%%%%%%%%%%%%
% TopoGuide -- ICRA-style manuscript (revised draft)
% Compile with ./build.sh
% Figures under pics/ are rendered when present; otherwise a labeled
% placeholder box is shown so the manuscript compiles standalone.
%%%%%%%%%%%%%%%%%%%%%%%%%%%%%%%%%%%%%%%%%%%%%%%%%%%%%%%%%%%%%%%%%%%%%%%%%%%%%%%%

\documentclass[letterpaper,10pt,conference]{ieeeconf}
\IEEEoverridecommandlockouts
\overrideIEEEmargins

\usepackage{amsmath}
\usepackage{amssymb}
\usepackage{booktabs}
\usepackage{graphicx}
\usepackage{xcolor}
\usepackage{placeins}
\usepackage{xstring}

\newcommand{\vect}[1]{\boldsymbol{#1}}
\newcommand{\mat}[1]{\mathbf{#1}}
\newcommand{\R}{\mathbb{R}}
\newcommand{\blank}{--}

% Render the figure when available; otherwise insert a labeled placeholder so
% that the manuscript compiles without the pics/ directory.
\newcommand{\safefig}[2][]{%
  \IfFileExists{#2}{\includegraphics[#1]{#2}}%
  {\begingroup\setlength{\fboxsep}{4pt}%
   \framebox[\dimexpr\linewidth-2\fboxsep-2\fboxrule]{%
     \parbox{0.88\linewidth}{\centering\vspace{1.6em}%
       \footnotesize\textit{[Figure placeholder]}\\[2pt]%
       \begingroup\footnotesize\ttfamily
       \edef\safefigpath{\detokenize{#2}}%
       \StrSubstitute{\safefigpath}{/}{/\allowbreak}[\safefigpathA]%
       \StrSubstitute{\safefigpathA}{_}{\_\allowbreak}[\safefigpathB]%
       \safefigpathB\endgroup%
       \vspace{1.6em}}}\endgroup}}

\title{\LARGE \bf
TopoGuide: A Plug-and-Play Topological Route Layer\\
for Language-Guided Aerial Navigation
}

% Author block anonymized for the double-anonymous review process.
\author{Anonymous Authors}

\begin{document}
\maketitle
\thispagestyle{empty}
\pagestyle{empty}

%%%%%%%%%%%%%%%%%%%%%%%%%%%%%%%%%%%%%%%%%%%%%%%%%%%%%%%%%%%%%%%%%%%%%%%%%%%%%%%%
\begin{abstract}
Local obstacle avoiders perform well in dense small-obstacle fields.
Their mapping modules, however, lack a unified representation for
obstacles much larger than the sensing horizon.  Sensor limits make
this gap worse: an obstacle can be visually obvious while depth
returns no measurement until the vehicle is too close to react.
Semantic heatmaps provide a beyond-visual-range cue that can guide
route selection early in a mission, before reliable geometry exists.
This paper presents TopoGuide, a plug-and-play topological route
layer that combines the sparse, persistent representation of a
verified topology graph with the long-range guidance of
open-vocabulary semantic heatmaps.  A lightweight FrontierGCN ranks
route directions directly on the online skeleton, while A* and
collision checks remain authoritative.  The accepted route is
exported only as a moving local goal, allowing existing
local planners to keep their fast small-obstacle avoidance while
gaining large-obstacle circumnavigation, without retraining or
modification of the planner core.  The semantic heatmap is
resolution-sensitive: higher input resolution carries richer scene
information.  The fusion interface supports various input resolutions.  In a
mission featuring a single 60-m-scale block that exceeds the depth
range, TopoGuide raises success from $0\%$ to $100\%$ for the
end-to-end planner YOPO, and to $70\%$ and $90\%$ for the
conventional pipelines EGO-Planner and SUPER.  A height-changing
mission with power-line obstacles and tree canopies further confirms
the same interface in volumetric flight.  TopoGuide acts as a
navigator: it remembers and selects the route, while an
interchangeable local planner executes only the next goal.
\end{abstract}

%%%%%%%%%%%%%%%%%%%%%%%%%%%%%%%%%%%%%%%%%%%%%%%%%%%%%%%%%%%%%%%%%%%%%%%%%%%%%%%%
\section{INTRODUCTION}
Aerial missions in forests, cities, and mines demand simultaneous
obstacle avoidance at vastly different scales, from branches and
wires to buildings and large structures.  Receding-horizon local
planners address the small-scale end of this spectrum well: fast
geometric planners and learned policies use depth and vehicle state
to generate short-horizon motion within milliseconds
\cite{rappids,yopo,loquercio}, and optimization-based pipelines
search and refine trajectories locally
\cite{egoplanner,fastplanner,super}.  Both families have demonstrated
effective avoidance in their evaluated dense environments.  Their
mapping modules, however, lack a unified representation that supports
both scales: local grids, distance fields, and lazy depth memories
remain confined to a sensing window \cite{voxblox}.  These planners are fast but carry no route memory: a single
observation cannot preserve an early left--right commitment around a
long obstacle, and the decision can be reversed once the fork leaves
the field of view.  A low-resolution global map \cite{voxblox} is a
natural remedy.  It can route around large structures.  Coarse
resolution, however, inflates narrow passages and closes them
entirely, making dense regions untraversable precisely where the
local planner excelled.  Sparse topological maps resolve this
asymmetry: they preserve verified free-space connectivity at local
geometric resolution while extending route memory far beyond the
sensing horizon \cite{yamauchi,sptm,ving,epic,bubble}.

\begin{figure}[!t]
\centering
\safefig[width=\columnwidth]{pics/teaser_candidate_v4.png}
\caption{Schematic of TopoGuide's layered route decision beyond the
depth-sensing horizon. (a)~The RGB image, semantic response, and depth map
are synchronized at graph time $t^*$; white depth pixels indicate no return
beyond $20$~m. (b)~FrontierGCN selects the central direction at the route
fork before the large obstacle, combines semantic observations with the
verified graph and safety checks, and provides the selected route as a
frontier goal. Blue-gray dots and links are verified graph nodes and edges;
red crosses are high-confidence far-field semantic observations, while gray
crosses are lower-confidence semantic candidates. The teal line is the
accepted A* topology path, the orange dashed line is its polynomial witness,
and the magenta diamond is the moving local goal.}
\label{fig:teaser}
\end{figure}


Onboard range sensors, including depth cameras and lightweight
LiDARs, are bounded in range, field of view, and angular resolution.
Thin or distant structures, such as power lines and building
facades, can be visually conspicuous yet return no reliable range
measurement until the vehicle is too close to react.  Open-vocabulary
semantic heatmaps offer a complementary beyond-visual-range cue
\cite{clip,openseg,pearl}: language-conditioned image responses
can mark hazardous regions early in a mission, before reliable
geometry exists.  Existing language-guided navigation systems exploit
this cue to decide \emph{where to go} \cite{vlfm,rayfronts}, but they
typically operate more slowly than the local controller and hand the
executor a goal rather than a collision-certified route.  They do not
address \emph{where not to go} at flight speed.

This paper presents TopoGuide, a plug-and-play topological route
layer that closes this gap.  TopoGuide maintains a sparse, persistent
graph of verified free space \cite{bubble,epic} and fuses it with the
long-range guidance of open-vocabulary semantic heatmaps
\cite{pearl,clip}.  A lightweight FrontierGCN ranks candidate route
directions directly on the online skeleton \cite{kipf2017gcn}, while
A* search and collision checking retain acceptance authority.  The
accepted route is exposed to the executor solely as a moving
local goal.  Existing local planners therefore retain
their small-obstacle agility and additionally acquire large-obstacle
circumnavigation, with no retraining or modification of the planner
core.  In this division of labor, TopoGuide acts as a navigator that
remembers and selects the route, while an interchangeable local
planner executes only the next goal \cite{ving}.

Experiments on a 140-m course containing both dense small obstacles
and a single 60-m-scale block show that three standalone executors,
an end-to-end reactive policy and two conventional pipelines, fail to
complete the mission, whereas the same planners guided by TopoGuide
achieve success rates of $100\%$, $70\%$, and $90\%$.  A
height-changing mission with power-line obstacles and tree canopies
further validates the interface in volumetric flight.

The main contributions are:
\begin{itemize}
  \setlength\itemsep{0pt}
  \item \textbf{A plug-and-play topological route layer} that unifies
  small- and large-scale obstacle avoidance: verified bubbles and
  witness edges retain route memory, exposed to planners only through
  a moving local goal (Secs.~\ref{sec:topology}
  and~\ref{sec:interface}).
  \item \textbf{A semantic heatmap fusion scheme} that supplies
  beyond-visual-range guidance: far-field frontiers and witness-level
  cost shape route preference without overriding geometric validity
  (Sec.~\ref{sec:semantics}).
  \item \textbf{A FrontierGCN frontier policy} that learns route
  direction selection on the online skeleton, trained on privileged
  A* labels from global environment maps, with recurrent temporal
  state smoothing unreliable semantic observations
  (Sec.~\ref{sec:gcn}).
  \item \textbf{A cross-executor evaluation} in which one route-layer
  implementation lifts three heterogeneous planners, end-to-end and
  conventional, from $0\%$ to up to $100\%$ success on a mission none
  completes alone (Sec.~\ref{sec:experiments}).
\end{itemize}

\begin{figure*}[!t]
\centering
\safefig[width=0.98\textwidth]{pics/system_architecture_topology_v3.pdf}
\vspace{-2pt}
\caption{TopoGuide as a plug-and-play route layer.  Depth and odometry maintain
persistent verified topology; RGB and text shape route cost, and a topology
GCN ranks frontier directions.  A* and collision checks accept a
witness before the moving-goal interface drives a goal-conditioned local
executor.  The executor receives only depth, state, and the local goal;
topology, semantics, and route memory remain upstream.}
\label{fig:system_architecture}
\end{figure*}

%%%%%%%%%%%%%%%%%%%%%%%%%%%%%%%%%%%%%%%%%%%%%%%%%%%%%%%%%%%%%%%%%%%%%%%%%%%%%%%%
\section{RELATED WORK}
\subsection{Local Planning and Topological Route Memory}
Local planners for aerial robots fall into two families.  Fast
geometric planners and learned policies use depth and vehicle state
to generate short-horizon motion commands or trajectories within
milliseconds \cite{rappids,yopo,loquercio}, whereas
optimization-based pipelines search and refine trajectories in local
maps \cite{egoplanner,fastplanner} or plan directly on point
clouds \cite{super}.  Both families have demonstrated effective
avoidance in their evaluated dense environments, but these methods
generally do not maintain a persistent, reusable route representation
across local replanning.

Sparse topological representations complement such planners with
route memory.  Frontier-based exploration retains frontier and
explored-space structure \cite{yamauchi}; bubble-based structures
support local replanning \cite{bubble}; and skeleton-based
topological graphs organize exploration in complex environments
\cite{epic}; and a persistent visibility graph enables
route recall and backtracking in partially known environments
\cite{farplanner}.  Learned visual memories and visual-goal
policies further show that retained experience benefits
longer-horizon navigation in ground and simulated domains
\cite{sptm,ving}.  We build on these components for goal-directed
flight: a topological skeleton with persistent, collision-checked
witness paths, shared by geometric search and the learned frontier
policy, and layered above an unchanged local planner core.

\subsection{Semantic Heatmaps in Navigation}
Open-vocabulary models align image regions with free-language queries
\cite{clip,openseg,pearl} and have been lifted into spatial
representations such as visual-language occupancy maps
\cite{vlmaps}.  For navigation, VLFM
builds frontier value maps from vision-language scores \cite{vlfm},
and RayFronts couples semantic ray frontiers with a dense onboard
occupancy map \cite{rayfronts}.  Their semantic update rates are
generally lower than those of the local controller, and they do not
directly provide collision-certified executable trajectories: in the
cited systems, semantic evidence is primarily used for frontier or
goal selection, while collision-certified route costs are handled
separately.  We
address the complementary question of where to avoid at flight speed:
heatmap responses become provisional far-field frontiers and
witness-level route cost.  This is consistent with risk-aware
planning \cite{blackmore}, where uncertain evidence modifies route
cost while geometric feasibility remains separately enforced.

\subsection{Graph Neural Policies}
Graph neural networks provide a general message-passing pattern that
aggregates relational state over sparse connectivity \cite{kipf2017gcn}.
In navigation, learned policies over topological or
belief graphs select long-range subgoals and frontier nodes for visual
navigation and exploration
\cite{chaplot2020ntslam,chen2020gnnexplore,cao2023ariadne}, and have
been validated on an aerial platform \cite{herrera2023gnnexplore}.
On the training side, global search has been distilled into learned
policies through imitation of heuristic search \cite{bhardwaj2017imitation},
differentiable A* search \cite{yonetani2021neuralastar}, and
privileged expert or critic supervision \cite{loquercio,cao2024largescale}.
We adopt this pattern for high-rate route selection on an online
route-memory skeleton, using edge-weighted convolution with global
topology pooling and a recurrent direction state; labels are computed
by A* on global environment maps, while the deployed policy observes
only the partial online skeleton.
%%%%%%%%%%%%%%%%%%%%%%%%%%%%%%%%%%%%%%%%%%%%%%%%%%%%%%%%%%%%%%%%%%%%%%%%%%%%%%%%
\section{TOPOLOGICAL MAPPING AND ROUTING}
\label{sec:method}
Let the vehicle state at time $t$ be
$\vect{x}_t=(\vect{p}_t,\vect{v}_t,\vect{a}_t,\vect{R}_t)$.
The sensors provide an RGB image $I_t$, a depth image $D_t$, and odometry.
The mission specifies a world-frame goal $\vect{g}$ and a text query $q$
describing regions to avoid.  No prior geometric map is available.

TopoGuide resolves the mission into a three-level goal hierarchy.  The
\emph{mission goal} $\vect{g}$ is the global task input and may lie far
beyond the sensing range.  The \emph{frontier goal}
$\vect{g}^{\mathrm{front}}_t$ is a route anchor selected on the topology
graph: it commits to a branch before that branch is fully observed.  The
\emph{local goal} $\vect{g}^{\mathrm{loc}}_t$ is sampled ahead on the
accepted witness, within the executor's depth envelope, and is the only
signal that crosses the planner boundary.  The planner must output a
dynamically feasible local trajectory $\tau_t$.  Formally, mapping, route
selection, and execution compose as
\begin{align}
  \mathcal{G}_t
    &= \Pi_{\mathrm{map}}(\mathcal{G}_{t-1},D_t,\vect{x}_t), \\
  (\mathcal{W}_t,\vect{g}^{\mathrm{front}}_t,
   \vect{g}^{\mathrm{loc}}_t,\hat c_t)
    &= \Pi_{\mathrm{route}}(\mathcal{G}_t,I_t,q,\vect{g}), \\
  \tau_t
    &= \Pi_{\mathrm{exec}}(D_t,\vect{o}_t).
\end{align}
Each layer has a deliberately narrow interface.  $\Pi_{\mathrm{map}}$
consumes depth, odometry, and its previous graph state; neither the semantic
model nor the local planner appears in its update.  $\Pi_{\mathrm{route}}$
produces the accepted witness $\mathcal{W}_t$ together with the frontier and
local goals and, when enabled, a frontier index $\hat c_t$ from the learned
policy; this index is a ranking hint, not a trajectory or safety decision.
For YOPO, the released policy receives its unmodified depth channel,
body-frame motion, and the local goal expressed in the body frame:
\begin{equation}
 \vect{o}_t = [\vect{v}^{B}_t,\vect{a}^{B}_t,
                (\vect{g}^{\mathrm{loc}}_t-\vect{p}_t)^{B}] \in \R^9.
\end{equation}
EGO-Planner and SUPER receive the same goal through a relay node.  No graph
tensor, semantic feature, or route representation crosses the executor
boundary; this division mirrors hierarchical topological control
\cite{ving}.

Bubble centers, checked witnesses, and moving goals are represented in
$\R^3$.  An optional projection onto one navigation height supports
controlled planar comparisons; with the projection disabled, graph
construction and route search preserve altitude, and the same witness
interface commands vertical as well as horizontal motion.

%%%%%%%%%%%%%%%%%%%%%%%%%%%%%%%%%%%%%%%%%%%%%%%%%%%%%%%%%%%%%%%%%%%%%%%%%%%%%%%%
\subsection{Topology Graph Construction and Maintenance}
\label{sec:topology}
The route layer maintains a sparse topological skeleton of verified free
space.  Collision-free bubbles are clustered into vertices
\begin{equation}
 v_i=(\vect{c}_i,\rho_i,z_i,h_i),
\end{equation}
where $\vect{c}_i$ is the center, $\rho_i$ the clearance, $z_i$ the semantic
state, and $h_i$ a persistent identity.  Each edge stores a checked witness
polyline $\pi_{ij}$ and its length $L_{ij}$; the witness, rather than the
chord between endpoints, is the geometric object used by both collision
checking and semantic fusion.

Construction runs asynchronously and independently of the semantic model,
the learned policy, and the local executor: depth frames are voxel-filtered
into obstacle points, a background worker rebuilds bubbles, regions, and
candidate edges, and a complete snapshot is atomically swapped into the
planner.  No occupancy grid or ESDF is retained.  Persistent identities
survive graph replacement, while stale bubbles are deleted after two
consecutive rebuild misses.  The minimum witness clearance $\bar\rho_{ij}$
contributes a soft preference
\begin{equation}
 C_{\mathrm{clr}}(e_{ij})=
 \lambda_{\mathrm{clr}}L_{ij}
 \left(\frac{\rho^\star}{\rho^\star+\bar\rho_{ij}}\right)^2,
\label{eq:clearance}
\end{equation}
but never turns a geometrically valid edge into an invalid one.

%%%%%%%%%%%%%%%%%%%%%%%%%%%%%%%%%%%%%%%%%%%%%%%%%%%%%%%%%%%%%%%%%%%%%%%%%%%%%%%%
\subsection{Semantic Heatmap Fusion}
\label{sec:semantics}
A frozen PEARL model maps $(I_t,q)$ to a patch heatmap.  The heatmap is
mean-pooled on a coarse grid, and each retained patch is projected at a
fixed optical-axis depth $\ell_s$:
\begin{equation}
 \vect{y}_{u,t}=\vect{p}_t+\mat{R}_t
 \left(\vect{t}_{\mathrm{cam}}+\ell_s\widetilde{\vect{d}}_u\right).
\label{eq:projection}
\end{equation}
Because the projection does not consult measured depth, an empty far-field
return cannot erase semantic evidence.  Each projection carries a
score--confidence tuple $(s_i,c_i)$ associated with a persistent graph
identity.

High-scoring, spatially separated projections become provisional semantic
frontiers.  A provisional frontier may attach only to nearby verified
backbone nodes, and its witness is checked against all currently known
obstacles; it cannot connect two verified corridors through unknown space.
When later geometry creates a bubble at the same position, the identity is
promoted to a verified node.  A frontier is therefore a route hypothesis,
never a free-space certificate.

Semantic evidence is fused with complete witness edges rather than only with
endpoints.  For semantic node $k$ and edge witness $\pi_{ij}$,
\begin{equation}
\begin{aligned}
 d_{k,ij}&=\min_{\vect{x}\in\pi_{ij}}
             \|\vect{y}_k-\vect{x}\|_2,\\
 r_{ij}&=\max\Big\{\tfrac12(s_i c_i+s_jc_j),\\[-2pt]
 &\qquad \max_{k:d_{k,ij}\le R_s}s_kc_k
 \exp[-d_{k,ij}^2/(2(R_s/2)^2)]\Big\}.
\end{aligned}
\label{eq:semantic_fusion}
\end{equation}
The resulting preference cost
\begin{equation}
 C_{\mathrm{sem}}(e_{ij})=
 \lambda_{\mathrm{sem}}L_{ij}
 [-\log(\max(\epsilon,1-r_{ij}))]
\label{eq:semantic_cost}
\end{equation}
reshapes route preference but never overrides geometric validity.

%%%%%%%%%%%%%%%%%%%%%%%%%%%%%%%%%%%%%%%%%%%%%%%%%%%%%%%%%%%%%%%%%%%%%%%%%%%%%%%%
\subsection{Route Search and Goal Interface}
\label{sec:interface}
Reachable vertices within radius $R_g$ are searched with the edge cost
$C(e)=W_{ij}+C_{\mathrm{clr}}(e)+C_{\mathrm{sem}}(e)$.  The selected witness
is stored in world coordinates and held until progress reaches a fixed
fraction or its geometric check fails; a fresh search then replaces the
complete route, and no graph-node pointers are retained across atomic
replacement.

The accepted witness, rather than a modified local policy, carries the route
decision.  Let $\mathcal{W}_t(s)$ be its arc-length parameterization and
$s_t$ the closest forward progress to the vehicle.  The local goal is
\begin{equation}
 \vect{g}^{\mathrm{loc}}_t=
 \mathcal{W}_t\!\left(\min(s_t+\ell_{\mathrm{loc}},L_{\mathcal W})\right),
 \qquad \ell_{\mathrm{loc}}=15~\mathrm{m}.
\label{eq:local_goal}
\end{equation}
The farther $\vect{g}^{\mathrm{front}}_t$ identifies the selected branch but
is not sent to the executor.  When a witness is replaced, the vehicle is
projected onto the new polyline before Eq.~\eqref{eq:local_goal} is
evaluated.

The executor contract is correspondingly narrow: a local planner must accept
a position goal and produce executable commands.  The original YOPO
checkpoint consumes the goal through its ordinary relative-goal input;
EGO-Planner and SUPER receive it through a relay node.  No adapter alters a
planner's search, optimization, or control law, and no corridor, bubble
sequence, or graph feature crosses the interface.  A stale or invalid
witness is rejected upstream and triggers a new search.

%%%%%%%%%%%%%%%%%%%%%%%%%%%%%%%%%%%%%%%%%%%%%%%%%%%%%%%%%%%%%%%%%%%%%%%%%%%%%%%%
\section{FRONTIERGCN DIRECTION POLICY}
\label{sec:gcn}
The FrontierGCN is TopoGuide's learned route-selection component.  It ranks
five yaw directions, left $40^\circ$ to right $40^\circ$ in $20^\circ$
steps, or, in the volumetric configuration, 15 yaw--pitch directions formed
with pitch rows of $+20^\circ$, $0^\circ$, and $-20^\circ$.  Its output is
only a frontier proposal: it neither creates a route nor issues an executor
command, leaving geometric acceptance to A*.

At each inference tick, skeleton nodes and edges form the input graph,
and virtual candidates are inserted $10$~m ahead of the vehicle at the
candidate directions and connected to their nearest skeleton node with
inverse-distance edge weights.  Each node carries a 20-dimensional feature
vector (24 in the volumetric variant) encoding normalized world- and
body-frame position, bubble clearance and degree, semantic score and
confidence, candidate identity, and mission-goal terms.  Two edge-weighted
graph-convolution layers \cite{kipf2017gcn} encode the graph, global mean
pooling supplies graph-level context, a GRU fuses the previously selected
direction, and an MLP scores the candidates.  A* then tests reachability and
collision validity exactly as it does for a hand-designed ranking.
Figure~\ref{fig:gcn_model} summarizes the deployed policy.  The volumetric
variant preserves three-dimensional graph coordinates and adds one
convolution layer.

\begin{figure}[t]
\centering
\makebox[\columnwidth][l]{\footnotesize\textbf{(a)} Deployed FrontierGCN}
\\[-1pt]
\safefig[width=0.98\columnwidth]{pics/gcn_model_architecture.pdf}
\vspace{2pt}
\makebox[\columnwidth][l]{\footnotesize\textbf{(b)} Privileged training-data generation}
\\[-1pt]
\safefig[width=0.98\columnwidth]{pics/gcn_training_data_log.pdf}
\vspace{-2pt}
\caption{FrontierGCN route-direction policy and its offline supervision.  (a)~Two
weighted graph-convolution layers encode the online skeleton and five virtual
candidates; pooled recurrent context produces a five-way route proposal, with
A* downstream.  (b)~Offline training-data construction follows the logged
point cloud and mission direction, an expert A* search for the ground-truth
path, and FrontierGCN's direction selection/output.  The graph is shown only
as a sparse top-down sketch; no semantic candidates are used in this label
generation.  The privileged A* path and direction target are unavailable at
deployment.}
\label{fig:gcn_model}
\end{figure}

Training labels are privileged.  For the planar policy, the complete
static ground-truth cloud is rasterized at $0.75$~m and inflated by the
$1.2$~m vehicle radius; eight-neighbor A* returns the direction of the
first path point at least $35$~m ahead.  The volumetric policy is labeled
by 26-connected A* in the inflated mesh volume.  Out-of-view and incomplete
paths are discarded, training minimizes class-weighted cross-entropy, and
session-level splits keep adjacent frames out of both the training and test
sets.  The ground-truth geometry is used only to
construct offline labels and is not available to the deployed policy.  The
resulting classifier is evaluated separately in
Sec.~\ref{sec:heuristic_eval}; its accuracy is not a safety claim.

%%%%%%%%%%%%%%%%%%%%%%%%%%%%%%%%%%%%%%%%%%%%%%%%%%%%%%%%%%%%%%%%%%%%%%%%%%%%%%%%
\section{EXPERIMENTS}
\label{sec:experiments}
The evaluation examines four complementary aspects of the system: closed-loop
compatibility of the topological goal interface with different local planners,
the effects of semantic cost and the GCN policy on route selection,
height-changing 3-D navigation, and the online graph-search workload under
persistent geometry.
Closed-loop and profiling trials are conducted in Unreal Engine 5.7 with
Colosseum, an AirSim-compatible simulator \cite{airsim}, using synchronized
RGB-D and odometry.  Unless noted otherwise, the primary experiments use a
1.6-m planning layer and one-way mission from $(0,0,1.6)$ to
$(0,140,1.6)$~m; neither the obstacle layout nor the nominal start and goal is
randomized across methods.  Section~\ref{sec:3d_eval} separately evaluates a
volumetric Map4 mission with the fixed-layer constraints disabled.

\subsection{Evaluation Protocol}
\label{sec:protocol}
\emph{Trial control and statistics.}  The simulation is reset before every
trial.  A goal is issued only after the vehicle remains within 0.10~m of the
common start, below 0.10~m/s, and within 0.10~rad of the mission heading for
1~s.  Startup failures are excluded from navigation outcomes, while collision
and timeout remain in the outcome rates; completion metrics use successful
trials only.

\emph{Metrics.}  Closed-loop metrics are success, collision, timeout, path
length $L$, flight time $T$, $\bar v=L/T$, and maximum speed.  For failed
flights, distance traveled before termination is reported separately as
$L_{\rm obs}$ and never mixed into successful-flight completion means.
Profiling separates point processing, graph construction, active A*, and the
complete planning tick.

\emph{Baseline fairness.}  EGO-Planner \cite{egoplanner} and SUPER
\cite{super} retain their planner cores but require simulator-specific cloud,
goal, and command adapters.  Their executor configurations and timeouts differ
from YOPO's, and recent benchmarking shows that such stacks are sensitive to
latency, field of view, and scenario difficulty \cite{flightbench}.  Their
failures therefore characterize these deployments, not universal planner
capability.  The plug-in claim concerns the narrow goal interface and absence
of planner-core changes, not strict experimental interchangeability.  The
no-graph rows are direct-goal executor controls rather than competitive global
planning baselines; Fig.~\ref{fig:integrated_results} therefore summarizes the
route-interface deployment, not superiority over map-based global A* or
sampling-based planners.

\subsection{Cross-Planner Closed-Loop Evaluation}
\subsubsection{Methods and Conditions}
The primary comparison pairs each standalone executor with a TopoGuide-guided
variant.  All TopoGuide variants run the same graph-node implementation,
semantic route formulation, A* acceptance, witness memory, and local-goal
interface.  YOPO consumes the goal directly; EGO-Planner and SUPER use the
adapter described in Sec.~\ref{sec:interface}.  Their internal planners do not
consume TopoGuide data structures.  Each row in the plug-in comparison contains
ten valid trials.  Table~\ref{tab:params} lists route-layer parameters shared
by the principal YOPO ablations.

\begin{table}[t]
\caption{Fixed Parameters Used by the Current Implementation}
\label{tab:params}
\centering
\small
\resizebox{\columnwidth}{!}{%
\begin{tabular}{lc@{\quad}lc}
\toprule
Parameter & Value & Parameter & Value \\
\midrule
Point voxel & 0.10 m & Point budget & 100,000 \\
Frame history & current frame only & Graph update & 100 ms \\
Graph rebuild & 100 ms & Search radius & 45 m \\
Local lookahead & 15 m & Goal-connect budget & 20 ms \\
$\lambda_{\rm sem}$ & 2.0 & $R_s$ & 8 m \\
$\lambda_{\rm clr}$ & 2.0 & $\rho^\star$ & 1.2 m \\
Score update & latest match & Edge-memory discount & disabled \\
Optical depth $\ell_s$ & 35 m & Node merge distance & 2.5 m \\
Semantic candidates/frame & $\leq 10$ & Virtual-node cap & 512 \\
Semantic candidates/edge & 8 & Semantic grid & $5\times3$ \\
\bottomrule
\end{tabular}
}
\end{table}

\subsubsection{Environment and Task}
All methods traverse the same ordered obstacle course
(Fig.~\ref{fig:integrated_results}).  Along the
$+y$ mission direction, the first approximately 25~m contain a dense field of
small obstacles.  The latter half is dominated by a single approximately
60-m-scale block, spanning roughly $y=68$--$122$~m and
$x=-27$--$26$~m, which forces an early commitment to a left or right bypass.
Its extent exceeds the 20-m depth clip and the local executor's immediate
reaction horizon, making persistent branch selection the central task.  The
ground-truth mesh spans $140.0\times132.0$~m.  In the evaluation
corridor $x\in[-45,45]$~m, $y\in[0,120]$~m, $3.91\%$ of 0.5-m cells intersect
ground-truth geometry in the navigation slab $z\in[0.6,2.6]$~m.  Ground truth
is used only for environment characterization and is not available to any method
online.  Matched language trials use the query ``blocks, wall.''

\subsubsection{Results}
Figure~\ref{fig:integrated_results} combines representative trajectory
overlays, closed-loop outcome rates, and online-resource time series.  Each
upper panel overlays one standalone executor with its TopoGuide-guided
counterpart in a common Map2 frame.  The aligned lower panel reports the
incremental route-layer wall time for that column's representative TopoGuide
run.

\begin{figure*}[!ht]
\centering
\safefig[width=0.98\textwidth]{pics/experiments/map2_0_140_1p6/figure4_integrated_topoguide.pdf}
\vspace{-2pt}
\caption{Integrated cross-executor evaluation.  The upper row overlays each
standalone executor (blue, orange, or green) with its TopoGuide counterpart
(red) on the same Map2 obstacle backdrop; colors identify methods and do not
encode speed, and upper-panel annotations report aggregate success.  The lower
row shows the online route-layer wall-time profile for the corresponding guided
run.  All curves are wall time in ms, binned over normalized mission time, and
trace point processing, graph construction, the complete planning tick, and
A* search latency.  Standalone logs do not expose comparable internal timing,
so the lower row isolates TopoGuide's added route-layer workload rather than
total executor compute.  Stars mark completed goals and crosses mark incomplete
baseline flights.}
\label{fig:integrated_results}
\end{figure*}

All three standalone configurations terminated without a successful mission,
whereas
the TopoGuide-guided YOPO, EGO-Planner, and SUPER variants completed 10/10,
7/10, and 9/10 trials.  This demonstrates that the same topological goal
contract is executable by three different planner cores.  Values marked
in the upper panels are aggregate ten-run outcomes, while the trajectories are
representative; failed-flight path and time are not treated as completion
metrics.  The TopoGuide (EGO)
deployment uses a 200-s mission timeout, while the other rows use 90~s.

The TopoGuide (YOPO) run produced no collision or timeout events: all ten trials
reached the goal and satisfied the final position and settling-speed
tolerances.  Relative to the 140-m straight-line mission, its mean path is
31.59~m longer, corresponding to a $22.6\%$ path overhead.  The longest trial
was trial 6 at $188.43$~m and $36.72$~s, which accounts for most of the path
length spread.  The mean terminal position error was $0.096\pm0.029$~m and
the mean terminal speed was $0.294\pm0.004$~m/s.

The latest executor regressions also expose distinct failure modes.  TopoGuide
with EGO-Planner had three collisions after 31.30, 83.41, and 132.15~m of
executed motion, spanning early, middle, and late portions of the mission.
TopoGuide with SUPER had one early collision after 6.95~m of executed motion
and no timeouts.  All ten trials in this evaluation were valid, with no
exclusions.  These terminal records distinguish collision from route
divergence or stall, but do not by themselves attribute a failure to the GCN,
accepted witness, or executor tracking.  Such attribution requires the
timestamped route and command trace rather than the terminal outcome alone.

\subsection{Route-Layer Ablations}
\subsubsection{Semantic Influence}
To test whether the route benefit depends on semantic evidence beyond the
depth envelope, we repeated the GCN stack with the same fixed mission,
90-s timeout, graph configuration, and query, while disabling PEARL,
semantic nodes, and semantic edge cost.  Completion metrics use successful
flights only, while $L_{\rm obs}$ records the path observed before collision
or timeout.  The resulting conventional flight metrics are reported in
Table~\ref{tab:gcn_semantic_ablation}.

\begin{table}[t]
\caption{Semantic route-layer ablation on the fixed mission with the GCN
heuristic fixed.
Semantic-off removes PEARL heatmaps, far-field nodes, and semantic edge cost;
geometry, A*, goal publication, and YOPO remain enabled.}
\label{tab:gcn_semantic_ablation}
\centering
\small
\resizebox{\columnwidth}{!}{%
\begin{tabular}{lcccccccc}
\toprule
Condition & Success (\%) & Collision (\%) & Timeout (\%) & $L_{\rm succ}$ (m) & $L_{\rm obs}$ (m) & Time (s) & Avg. speed (m/s) & Max. speed (m/s) \\
\midrule
Matched query & 100 & 0 & 0 & $171.59\pm6.57$ & \blank & $33.61\pm1.25$ & $5.105\pm0.045$ & $6.254\pm0.138$ \\
Semantic disabled & 60 & 40 & 0 & $169.47\pm3.16$ & $11.69\pm1.60$ & $32.34\pm2.70$ & $5.261\pm0.315$ & $6.404\pm0.158$ \\
Irrelevant query & 80 & 20 & 0 & $175.76\pm6.09$ & $20.32\pm2.67$ & $35.02\pm1.65$ & $5.023\pm0.127$ & $6.398\pm0.274$ \\
\bottomrule
\end{tabular}
}
\end{table}

Semantic influence changes the outcome rate without changing the geometric
validity contract: the semantic-on condition completed all ten flights,
whereas the semantic-off condition had four early collisions near
$y\approx10$~m.  The
successful-flight path and speed differences are secondary to this outcome
gap.  Together with the
depth-clipped projection visualized in Fig.~\ref{fig:teaser}, this behavior
is consistent
with semantic evidence affecting a route before the corresponding obstacle is
measured by depth.
The irrelevant-query condition (prompt ``trees'') completes 8/10 trials,
confirming that route selection responds to the alignment between the language
query and visual scene content.

\subsubsection{GCN Frontier-Policy Ablation}
To isolate the learned route-selection component, we compare TopoGuide's
hand-designed frontier ranking with the GCN policy while
holding the fixed mission, semantic query, route-cost parameters, and YOPO
executor fixed.

\begin{table}[t]
\caption{Frontier-Policy Ablation Within TopoGuide}
\label{tab:single_run_validation}
\centering
\small
\resizebox{\columnwidth}{!}{%
\begin{tabular}{lccccccc}
\toprule
Method & Success (\%) & Collision (\%) & Timeout (\%) & Path (m) & Time (s) & Avg. speed (m/s) & Max. speed (m/s) \\
\midrule
Hand-designed ($n=27$) & 88.9 & 11.1 & 0 & $199.14\pm14.12$ & $38.54\pm3.06$ & $5.172\pm0.097$ & $6.409\pm0.184$ \\
GCN ($n=10$) & 100.0 & 0.0 & 0.0 & $171.59\pm6.57$ & $33.61\pm1.25$ & $5.105\pm0.045$ & $6.254\pm0.138$ \\
\bottomrule
\end{tabular}
}
\end{table}

The GCN policy completed 10/10 trials, compared with 24/27 for the
hand-designed ranking, while reducing mean path length from 199.14~m to
171.59~m and mean flight time from 38.54~s to 33.61~s.  Only the
frontier-ranking policy is replaced; A*, witness acceptance, semantic route
cost, and the local-goal interface remain unchanged.

\subsection{GCN Policy Evaluation}
\label{sec:heuristic_eval}
The privileged dataset contains 5,337 frames, split by session into 3,844
training, 786 validation, and 707 held-out test frames.  A validation sweep
selects no additional switching penalty.  On
the untouched test split, the deployed checkpoint obtains $89.3\%$ exact and
$89.4\%$ macro accuracy; $99.4\%$ of predictions are within one adjacent
direction column, with a mean absolute column error of $0.115$.  The logged
hand-designed frontier policy matches $22.2\%$ of the strict 35-m labels.
These results measure route-direction prediction, not online safety.

\begin{table}[t]
\caption{GCN Direction Classification}
\label{tab:gcn_offline}
\centering
\small
\begin{tabular}{lc}
\toprule
Metric & Test \\
\midrule
Exact accuracy & 89.3\% \\
Macro accuracy & 89.4\% \\
Within one direction column & 99.4\% \\
Mean absolute column error & 0.115 \\
\bottomrule
\end{tabular}
\end{table}

Figure~\ref{fig:gcn_real_views} shows synchronized frames from the online
interface.  The RGB image, PEARL heatmap, selected GCN column, and mission-goal
direction are retained from the recorded experiment.

\begin{figure*}[t]
\centering
\begin{minipage}[t]{0.235\textwidth}
\centering
\safefig[width=\linewidth]{pics/0903_graph_semantic_frontier_building1.png}\\[-1pt]
{\footnotesize (a) Building view 1}
\end{minipage}\hfill
\begin{minipage}[t]{0.235\textwidth}
\centering
\safefig[width=\linewidth]{pics/0903_graph_semantic_frontier_building2.png}\\[-1pt]
{\footnotesize (b) Building view 2}
\end{minipage}\hfill
\begin{minipage}[t]{0.235\textwidth}
\centering
\safefig[width=\linewidth]{pics/0903_graph_semantic_frontier_forest1.png}\\[-1pt]
{\footnotesize (c) Forest view 1}
\end{minipage}\hfill
\begin{minipage}[t]{0.235\textwidth}
\centering
\safefig[width=\linewidth]{pics/0903_graph_semantic_frontier_forest2.png}\\[-1pt]
{\footnotesize (d) Forest view 2}
\end{minipage}
\caption{Recorded GCN-interface views from closed-loop flights.  Each
panel preserves the RGB frame and PEARL heatmap; red marks the proposed
frontier column and the
green dashed arrow marks the 35-m mission-goal direction.}
\label{fig:gcn_real_views}
\end{figure*}

\subsection{Height-Changing 3-D Navigation}
\label{sec:3d_eval}
We further evaluate the volumetric mode in a separate Map4 line-obstacle
mission.  The start and goal are $(0,0,3.5)$ and $(0,140,3.5)$~m, placing the
nominal flight height through the power-line corridor.  Both the graph-layer
projection and YOPO's fixed-altitude constraint are disabled.  A dedicated
15-direction checkpoint, trained on 2,200 graph states labeled by
26-connected A* in the privileged 3-D mesh volume, ranks three pitch rows and
five yaw columns; the online planner still observes only RGB-D, odometry, and
the text query ``line.''

\begin{table}[t]
\caption{Height-Changing Map4 Line-Obstacle Mission ($n=10$).  Continuous
values are computed over the successful flights.}
\label{tab:3d_line_navigation}
\centering
\small
\resizebox{\columnwidth}{!}{%
\begin{tabular}{lccccccc}
\toprule
Method & Success (\%) & Collision (\%) & Timeout (\%) & Path (m) & Time (s) & Avg. speed (m/s) & $z$ span (m) \\
\midrule
TopoGuide (GCN-3D) & 80 & 20 & 0 & $167.40\pm13.83$ & $36.98\pm5.30$ & $4.564\pm0.327$ & $2.225\pm0.477$ \\
\bottomrule
\end{tabular}
}
\end{table}

TopoGuide completes $80\%$ of the trials without timeout.  Every successful
flight changes altitude substantially: the per-flight vertical span ranges
from 1.731 to 2.934~m.  The shortest successful route is 149.203~m in
30.934~s and spans 1.808~m vertically.  The two failures are early collisions
after 10.907 and 10.019~m.  Together with the non-planar graph and executor
settings, these trajectories verify that the same topology--witness--goal
pipeline supports height-changing route execution rather than only planar
branch selection.

\subsection{Online Planning Workload}
The online search is restricted to a local working set even when the verified
topology persists outside the current sensing window.  The current
configuration keeps at most 100,000 point samples, rebuilds the skeleton every
100~ms, searches within 45~m, and caps virtual semantic nodes at 512.
The lower row of Fig.~\ref{fig:integrated_results} shows time-binned wall-time
traces for the representative TopoGuide run paired with each executor.  The
persistent mode uses 214/249 local
nodes at P50/P95 and expands 185/215 nodes per active A* search, compared with
197/231 and 173/194 when persistence is disabled.  Thus retaining the explored
backbone does not make every retained node part of each planning tick.  Over
normalized session time, the observed local-node P95 remains below 279 with
persistence and 248 without it.  At approximately 500 nodes and 3,200 directed
edges, GCN inference on an NVIDIA GeForce RTX 3090 adds 0.83~ms mean latency
to the online workload, with 0.60~ms at P50 and 1.97~ms at P95.
The traces separate point processing, graph construction, the complete
planning tick, and active A*.  Because the standalone logs lack the same
internal timing instrumentation, these panels characterize the incremental
TopoGuide layer and are not a comparison of total executor compute.

A route holds its accepted witness until a progress gate is reached or a
geometric check fails; this makes route switching observable and prevents a
partially rebuilt graph from leaking into the controller.  The semantic stream
runs at roughly 2~Hz, whereas graph and control updates remain asynchronous.

For reproducibility, each trial records the graph snapshot timestamp, selected
frontier column, witness clearance, A* replacement reason, local-goal distance,
and terminal outcome.  These fields let us distinguish a direction proposal
from the route actually accepted after collision checking.  They also expose
the failure sequence seen in the evaluation: a stale or low-clearance
witness is rejected, a fresh A* search is requested, and a moving vehicle may
reach the executor's start-failure condition before the replacement goal is
established.  This logging boundary is important because it attributes safety
events to the graph--planner--executor interface rather than to an offline GCN
classification score alone.

%%%%%%%%%%%%%%%%%%%%%%%%%%%%%%%%%%%%%%%%%%%%%%%%%%%%%%%%%%%%%%%%%%%%%%%%%%%%%%%%
\section{CONCLUSION AND DISCUSSION}
\label{sec:discussion}
TopoGuide is a topological route layer rather than another local planner:
it turns depth and odometry into persistent witness-verified structure,
combines language-conditioned route preference with a learned frontier
policy, and exposes only a moving local goal.  The same implementation
drove three heterogeneous executors without retraining, supporting the
central claim: long-horizon memory and language-conditioned route choice
can be added above interchangeable short-horizon executors.

Several limitations remain.  ``Plug-and-play'' means no planner retraining
once the goal and command interfaces are adapted; it does not remove
integration code.  The cross-executor study fixes one altitude layer and
one mission layout, and the volumetric evaluation covers a single
line-obstacle mission; randomized layouts, dynamic obstacles, and real
sensors remain untested.  Semantic projections remain hypotheses checked
only against known obstacles, and offline GCN accuracy measures direction
prediction, not online safety.  Future graph policies can incorporate
clearance and braking margin directly and mask invalid directions before
learned ranking.


\begin{thebibliography}{99}

\bibitem{rappids}
N. Bucki, J. Lee, and M. W. Mueller,
``Rectangular Pyramid Partitioning Using Integrated Depth Sensors (RAPPIDS):
A Fast Planner for Multicopter Navigation,'' \emph{IEEE Robot. Autom. Lett.},
vol. 5, no. 3, pp. 4626--4633, 2020.
\bibitem{yopo}
J. Lu, X. Zhang, H. Shen, L. Xu, and B. Tian,
``You Only Plan Once: A Learning-Based One-Stage Planner With Guidance
Learning,'' \emph{IEEE Robot. Autom. Lett.}, vol. 9, no. 7,
pp. 6083--6090, 2024.
\bibitem{loquercio}
A. Loquercio, E. Kaufmann, R. Ranftl, M. M\"uller, V. Koltun, and
D. Scaramuzza,
``Learning High-Speed Flight in the Wild,''
\emph{Sci. Robot.}, vol. 6, no. 59, 2021.
\bibitem{egoplanner}
X. Zhou, Z. Wang, H. Ye, C. Xu, and F. Gao,
``EGO-Planner: An ESDF-Free Gradient-Based Local Planner for Quadrotors,''
\emph{IEEE Robot. Autom. Lett.}, vol. 6, no. 2, pp. 478--485, 2021.
\bibitem{fastplanner}
B. Zhou, F. Gao, L. Wang, C. Liu, and S. Shen,
``Robust and Efficient Quadrotor Trajectory Generation for Fast
Autonomous Flight,'' \emph{IEEE Robot. Autom. Lett.}, vol. 4, no. 4,
pp. 3529--3536, 2019.
\bibitem{super}
Y. Ren, F. Zhu, G. Lu, Y. Cai, L. Yin, F. Kong, J. Lin, N. Chen, and
F. Zhang,
``Safety-Assured High-Speed Navigation for MAVs,''
\emph{Sci. Robot.}, vol. 10, no. 98, p. eado6187, 2025.
\bibitem{voxblox}
H. Oleynikova, Z. Taylor, M. Fehr, R. Siegwart, and J. Nieto,
``Voxblox: Incremental 3D Euclidean Signed Distance Fields for On-Board MAV
Planning,'' \emph{IEEE/RSJ Int. Conf. Intell. Robots Syst. (IROS)}, 2017.
\bibitem{yamauchi}
B. Yamauchi,
``A Frontier-Based Approach for Autonomous Exploration,'' in \emph{Proc.
IEEE Int. Symp. Comput. Intell. Robot. Autom.}, 1997, pp. 146--151.
\bibitem{sptm}
N. Savinov, A. Dosovitskiy, and V. Koltun,
``Semi-Parametric Topological Memory for Navigation,'' in \emph{Proc. Int.
Conf. Learn. Represent.}, 2018.
\bibitem{ving}
D. Shah, B. Eysenbach, G. Kahn, N. Rhinehart, and S. Levine,
``ViNG: Learning Open-World Navigation With Visual Goals,'' in \emph{Proc.
IEEE Int. Conf. Robot. Autom.}, 2021, pp. 13215--13222.
\bibitem{epic}
S. Geng, Z. Ning, F. Zhang, and B. Zhou,
``EPIC: A Lightweight LiDAR-Based AAV Exploration Framework for Large-Scale
Scenarios,'' \emph{IEEE Robot. Autom. Lett.}, vol. 10, no. 5,
pp. 5090--5097, 2025.
\bibitem{bubble}
Y. Ren, F. Zhu, W. Liu, Z. Wang, Y. Lin, F. Gao, and F. Zhang,
``Bubble Planner: Planning High-Speed Smooth Quadrotor Trajectories Using
Receding Corridors,'' in \emph{Proc. IEEE/RSJ Int. Conf. Intell. Robots
Syst.}, 2022, pp. 6332--6339.
\bibitem{clip}
A. Radford et al.,
``Learning Transferable Visual Models From Natural Language Supervision,''
in \emph{Proc. Int. Conf. Mach. Learn.}, 2021, pp. 8748--8763.
\bibitem{openseg}
G. Ghiasi, X. Gu, Y. Cui, and T.-Y. Lin,
``Scaling Open-Vocabulary Image Segmentation With Image-Level Labels,'' in
\emph{Proc. Eur. Conf. Comput. Vis.}, 2022.
\bibitem{pearl}
G. Pei, X. Jiang, X. Cai, T. Chen, Y. Yao, and B. Jeon,
``PEARL: Geometry Aligns Semantics for Training-Free Open-Vocabulary Semantic
Segmentation,'' in \emph{Proc. IEEE/CVF Conf. Comput. Vis. Pattern Recognit.},
2026, pp. 17927--17937.
\bibitem{vlfm}
N. Yokoyama, S. Ha, D. Batra, J. Wang, and B. Bucher,
``VLFM: Vision-Language Frontier Maps for Zero-Shot Semantic Navigation,''
in \emph{Proc. IEEE Int. Conf. Robot. Autom.}, 2024, pp. 42--48.
\bibitem{rayfronts}
O. Alama, A. Bhattacharya, H. He, S. Kim, Y. Qiu, W. Wang, C. Ho, N. Keetha,
and S. Scherer,
``RayFronts: Open-Set Semantic Ray Frontiers for Online Scene Understanding
and Exploration,'' in \emph{Proc. IEEE/RSJ Int. Conf. Intell. Robots Syst.},
2025, pp. 5930--5937.
\bibitem{kipf2017gcn}
T. N. Kipf and M. Welling,
``Semi-Supervised Classification with Graph Convolutional Networks,''
in \emph{Proc. Int. Conf. Learn. Represent.}, 2017.
\bibitem{farplanner}
F. Yang, C. Cao, H. Zhu, J. Oh, and J. Zhang,
``FAR Planner: Fast, Attemptable Route Planner Using Dynamic Visibility
Update,'' in \emph{Proc. IEEE/RSJ Int. Conf. Intell. Robots Syst.
(IROS)}, 2022, pp. 9--16.
\bibitem{vlmaps}
C. Huang, O. Mees, A. Zeng, and W. Burgard,
``Visual Language Maps for Robot Navigation,'' in \emph{Proc. IEEE Int.
Conf. Robot. Autom.}, 2023, pp. 10608--10615.
\bibitem{blackmore}
L. Blackmore, M. Ono, and B. C. Williams,
``Chance-Constrained Optimal Path Planning With Obstacles,'' \emph{IEEE
Trans. Robot.}, vol. 27, no. 6, pp. 1080--1094, 2011.
\bibitem{chaplot2020ntslam}
D. S. Chaplot, R. Salakhutdinov, A. Gupta, and S. Gupta,
``Neural Topological SLAM for Visual Navigation,'' in \emph{Proc.
IEEE/CVF Conf. Comput. Vis. Pattern Recognit.}, 2020.
\bibitem{chen2020gnnexplore}
F. Chen, J. D. Martin, Y. Huang, J. Wang, and B. Englot,
``Autonomous Exploration Under Uncertainty via Deep Reinforcement
Learning on Graphs,'' in \emph{Proc. IEEE/RSJ Int. Conf. Intell.
Robots Syst. (IROS)}, 2020, pp. 6140--6147.
\bibitem{cao2023ariadne}
Y. Cao, T. Hou, Y. Wang, X. Yi, and G. Sartoretti,
``ARiADNE: A Reinforcement Learning Approach Using Attention-Based
Deep Networks for Exploration,'' in \emph{Proc. IEEE Int. Conf.
Robot. Autom.}, 2023, pp. 10219--10225.
\bibitem{herrera2023gnnexplore}
E. P. Herrera-Alarc\'on, G. Baris, M. Satler, C. A. Avizzano, and
G. Loianno,
``Learning Heuristics for Efficient Environment Exploration Using
Graph Neural Networks,'' in \emph{Proc. Int. Conf. Adv. Robot.},
2023, pp. 86--93.
\bibitem{bhardwaj2017imitation}
M. Bhardwaj, S. Choudhury, and S. Scherer,
``Learning Heuristic Search via Imitation,'' in \emph{Proc. Conf.
Robot Learn. (CoRL)}, ser. PMLR, vol. 78, 2017, pp. 271--280.
\bibitem{yonetani2021neuralastar}
R. Yonetani, T. Taniai, M. Barekatain, M. Nishimura, and A. Kanezaki,
``Path Planning Using Neural A* Search,'' in \emph{Proc. Int. Conf.
Mach. Learn. (ICML)}, ser. PMLR, vol. 139, 2021, pp. 12029--12039.
\bibitem{cao2024largescale}
Y. Cao, R. Zhao, Y. Wang, B. Xiang, and G. Sartoretti,
``Deep Reinforcement Learning-Based Large-Scale Robot Exploration,''
\emph{IEEE Robot. Autom. Lett.}, vol. 9, no. 5, pp. 4631--4638, 2024.
\bibitem{airsim}
S. Shah, D. Dey, C. Lovett, and A. Kapoor,
``AirSim: High-Fidelity Visual and Physical Simulation for Autonomous
Vehicles,'' in \emph{Field and Service Robotics}, 2018, pp. 621--635.
\bibitem{flightbench}
S.-A. Yu, C. Yu, F. Gao, Y. Wu, and Y. Wang,
``FlightBench: Benchmarking Learning-Based Methods for Ego-Vision-Based
Quadrotors Navigation,'' in \emph{Proc. IEEE Int. Conf. Robot. Autom.},
2025.

\end{thebibliography}

\end{document}
